% Source: https://artofproblemsolving.com/community/c3h2262172_proofs_without_words_ii_by_roger_b_nelson
% Stable identity: aops-topic-2262172-q01
\textbf{(a)} Multiply by $4$ and subtract:
\[4S_n-S_n=\left(1+\dfrac14+\cdots+\dfrac1{4^{n-1}}\right)-\left(\dfrac14+\cdots+\dfrac1{4^n}\right)=1-\dfrac1{4^n},\]
so $S_n=\dfrac13\left(1-\dfrac1{4^n}\right)$. As $n$ grows, $\dfrac1{4^n}\to0$, so the partial sums approach $\dfrac13$, which is what the infinite sum means:
\[\dfrac14+\dfrac1{4^2}+\dfrac1{4^3}+\cdots=\boxed{\dfrac13}.\]

\textbf{(b)}

\begin{asy}
usepackage("amsmath");
size(4.5cm);
defaultpen(fontsize(10pt));
real s=4;
void sq(pair o,real w,int lev){
  fill(box(o,o+(w/2,w/2)),rgb(0.45,0.65,1));
  draw(box(o,o+(w,w)));
  draw((o.x+w/2,o.y)--(o.x+w/2,o.y+w)); draw((o.x,o.y+w/2)--(o.x+w,o.y+w/2));
  if(lev>0) sq(o+(w/2,w/2),w/2,lev-1);
}
sq((0,0),s,5);
draw(box((0,0),(s,s)),linewidth(1));
label("$\dfrac14$",(s/4,s/4)); label("$\dfrac1{16}$",(5s/8,5s/8),fontsize(8pt));
\end{asy}

At the first stage, the square minus its top-right quarter is an L-shape made of three equal quarter-squares, one of which is shaded: exactly $\dfrac13$ of the L-shape is shaded. The same happens inside the top-right quarter at the next stage, and so on. These L-shapes fill the whole square except the ever smaller top-right corner square, whose area $\dfrac1{4^n}$ shrinks to $0$. Since every L-shape is one-third shaded, the shaded area, which is $\dfrac14+\dfrac1{4^2}+\cdots$, is $\dfrac13$ of the area of the square, i.e. $\dfrac13$.
